% Shared preamble for TMUA worked solutions. Local edits here are overwritten.

% Reproducible PDFs: no timestamps, no trailer id, no build-path details.
\ifdefined\pdfsuppressptexinfo
  \pdfsuppressptexinfo=-1
  \pdfinfoomitdate=1
  \pdftrailerid{}
\fi

\usepackage{graphicx}
\usepackage{mathtools}
\usepackage{amssymb}
\usepackage{amsmath}
\usepackage{amsthm}
\usepackage{enumitem}
\usepackage{microtype}
\usepackage[table]{xcolor}
\usepackage{array}
\usepackage{booktabs}
\usepackage{tabularx}
\usepackage{longtable}
\usepackage{ragged2e}
\usepackage{fancyhdr}
\usepackage{etoolbox}
\usepackage[
  a4paper,
  left=0.8in,
  right=1in,
  top=1in,
  bottom=0.5in,
  headheight=14pt,
  headsep=0.4in
]{geometry}

\usepackage{pgfplots}
\pgfplotsset{compat=1.18}
\usepgfplotslibrary{fillbetween, polar}
\usetikzlibrary{calc, positioning}
\usetikzlibrary{angles, quotes}
\usetikzlibrary{arrows.meta}
\usetikzlibrary{decorations.pathreplacing, decorations.markings}
\usetikzlibrary{patterns}
\usetikzlibrary{intersections}

\usepackage{tocloft}
\usepackage[hidelinks]{hyperref}
\hypersetup{pdfauthor={danieljsmith.org}, pdfcreator={}, pdfsubject={}, pdfkeywords={}}
\urlstyle{same}
\tocloftpagestyle{djssolutions}
\setlength{\cftbeforesubsecskip}{2pt}
\setlength{\cftsubsecindent}{1.5em}
\setlength{\cftsubsecnumwidth}{0pt}
\renewcommand{\cftsubsecfont}{\small\color{black}}
\renewcommand{\cftsubsecpagefont}{\small\color{black}}
\renewcommand{\cftsubsecleader}{\hfill}

% `most` carries the breakable and enhanced libraries the solution box needs.
\usepackage[most]{tcolorbox}
\definecolor{scarlet}{RGB}{205, 40, 75}

\setlength{\parindent}{0pt}

% \djsSolutionsSetup{<paper title>}, called once after this file.
\newcommand{\djsSolnPaperTitle}{}
\newcommand{\djsSolutionsSetup}[1]{%
  \hypersetup{pdftitle={#1 Solutions}}%
  \renewcommand{\djsSolnPaperTitle}{#1}}

% \djsTitleSetup{<booklet>}{<title>}{<paper name>}{<questions>}{<minutes>},
% called once after the setup line, sets the title block that opens the
% contents page, as on the question paper: TMUA and the booklet in spaced
% burgundy capitals, the title, a hairline led by the solution box's short
% scarlet rule, then the paper's name with the question count and, when
% given, the time allowed. Its colours are the site's accent, rule and muted
% text colours.
\definecolor{djsBurgundy}{RGB}{118, 59, 67}
\definecolor{djsHairline}{RGB}{216, 213, 205}
\definecolor{djsSlash}{RGB}{170, 166, 158}
\definecolor{djsMuted}{RGB}{96, 101, 106}
\makeatletter
\newcommand{\djs@booklet}{}
\newcommand{\djs@title}{}
\newcommand{\djs@papername}{}
\newcommand{\djs@details}{}
\newcommand{\djs@caps}[1]{\textls[180]{\MakeUppercase{#1}}}
\newcommand{\djs@dot}{\hspace{0.45em}\textperiodcentered\hspace{0.45em}}
\newcommand{\djsTitleSetup}[5]{%
  \renewcommand{\djs@booklet}{#1}%
  \renewcommand{\djs@title}{#2}%
  \renewcommand{\djs@papername}{#3}%
  \renewcommand{\djs@details}{%
    #4~question\ifnum#4=1 \else s\fi
    \ifblank{#5}{}{\djs@dot#5~minutes}}}
\newcommand{\djsContentsTitle}{%
  {\raggedright\color{djsBurgundy}\footnotesize\bfseries
    \djs@caps{TMUA}{\color{djsSlash}\hspace{0.55em}/\hspace{0.55em}}%
    \djs@caps{\djs@booklet}\par}%
  \vspace{7pt}%
  {\raggedright\fontsize{24.88}{28}\selectfont\djs@title\par}%
  \vspace{9pt}%
  \noindent\makebox[0pt][l]{\color{djsHairline}\rule[0.9pt]{\linewidth}{0.6pt}}%
  {\color{scarlet}\rule{3.5cm}{2.2pt}}\par
  \vspace{5pt}%
  {\small\color{djsMuted}\djs@papername\hfill\djs@details\par}%
  \vspace{8pt}}
\makeatother

% The header appears only on the contents page and each question's opening
% page. Continuation pages use the empty default style.
\fancypagestyle{djssolutions}{%
  \fancyhead{}%
  \fancyfoot{}%
  \renewcommand{\headrulewidth}{0.9pt}%
  \fancyhead[L]{\textit{Solutions} - \djsSolnPaperTitle}%
  \fancyhead[R]{\href{https://danieljsmith.org/?utm_source=pdf&utm_campaign=tmua&utm_content=solutions}{danieljsmith.org}}%
}
\pagestyle{empty}

% \djsLastUpdated{<date>}, called once after the setup line: one line
% at the foot of the first page only, left-aligned under the left header at
% the 0.8in left margin. It is drawn over the page rather than set as a
% footer, so no page style changes.
\newcommand{\djsLastUpdated}[1]{%
  \AtBeginDocument{\AddToHookNext{shipout/foreground}{%
    \put(0.8in,-\dimexpr\paperheight-0.32in\relax){\makebox[0pt][l]{%
      \footnotesize Last updated #1}}}}}

\newcolumntype{Y}{>{\RaggedRight\arraybackslash}X}

% Question layout, identical to the question paper's.
\newcommand{\djsTocTopic}[1]{%
  \ifblank{#1}{}{%
    \texorpdfstring
      {\hspace{0.9em}{\footnotesize\itshape\color{black!45}#1}}
      { \textendash\ #1}}}
\newcommand{\djsTopicLine}[1]{%
  \ifblank{#1}{}{%
    {\raggedleft\footnotesize\itshape\color{black!45}#1\par}%
    \vspace{-0.6\baselineskip}}}
\newcommand{\questionitem}[1][]{%
  \item\phantomsection
  \addcontentsline{toc}{subsection}{%
    Question \number\value{enumi}\protect\djsTocTopic{#1}}%
}
\renewcommand{\contentsname}{Questions}
\newcommand{\djsFrontMatter}{%
  \vspace*{6pt}%
  \djsContentsTitle
  \vspace{14pt}%
  \tableofcontents
  \newpage
}

% The solution body: a white breakable box with a left scarlet rule. The
% rule continues onto a continuation page and stops where the solution
% stops, then turns into a short horizontal of the same weight so the
% solution visibly ends. The argument is the answer label.
\newcommand{\djsSolutionAnswer}[1]{%
  {\color{scarlet}\bfseries Answer\ \ #1}}
\newtcolorbox{djsSolution}[1]{%
  enhanced, breakable, frame hidden, sharp corners, colback=white,
  height fixed for=first and middle, height fill,
  extras unbroken and last={natural height},
  borderline west={2.2pt}{0pt}{scarlet},
  left=11pt, right=0pt, top=5pt, bottom=13pt,
  before skip=13pt, after skip=4pt,
  overlay unbroken and last={%
    \fill[scarlet] (frame.south west)
      rectangle ([xshift=3.5cm, yshift=2.2pt]frame.south west);},
  before upper={\setlength{\parskip}{6pt}%
                {\color{scarlet}\bfseries Solution}%
                \hfill\djsSolutionAnswer{#1}\par\vspace{4pt}}}

% Solution figures sit inside the box and are drawn smaller than question
% figures; \scalebox shrinks node text with the drawing.
\newcommand{\SolnFigureScale}{0.85}
\newsavebox{\djsSolnFigureBox}
\newenvironment{djsSolutionFigure}
  {\par\vspace{4pt}\begin{center}\begin{lrbox}{\djsSolnFigureBox}}
  {\end{lrbox}\scalebox{\SolnFigureScale}{\usebox{\djsSolnFigureBox}}%
   \end{center}\vspace{2pt}\par}
